\chapter{ANTECEDENTES}
\label{cap:antecedentes}

% ===========================================================================
% ESTADO DEL ARTE. Regla que separa este capítulo del Cap. 4:
%   - ANTECEDENTES  -> QUÉ ES y QUÉ EXISTE (conceptos, evolución, trabajos previos)
%   - MAT. Y MÉTODOS -> QUÉ ELEGÍ YO y POR QUÉ (alternativas, comparativas, ADR)
% Aquí NO va ninguna decisión del proyecto: aquí se prepara al lector para
% entenderla. Si una frase empieza por "se ha elegido", está en el capítulo
% equivocado.
% Cubre IT-20 (estado del arte y trabajos relacionados).
% ===========================================================================

\section{Introducción}

% Un párrafo: qué se revisa en este capítulo y en qué orden. Sirve de mapa.

\section{La orientación académica preuniversitaria}

\subsection{El problema de la elección de titulación}

% El punto de partida humano del TFG: qué se juega un estudiante de secundaria
% al elegir grado, y por qué la información oficial no le llega. Datos si los
% encuentras (abandono en primer curso, cambios de titulación).

\subsection{Recursos de orientación existentes}

% Qué tiene hoy a su disposición: webs oficiales, guías en PDF, jornadas de
% puertas abiertas, comparadores de titulaciones, orientadores de instituto.

\subsection{Limitaciones de los enfoques actuales}

% Por qué no bastan: exigen saber el término técnico exacto, no responden dudas
% abiertas, la información está dispersa. Esta subsección es la que abre hueco
% a tu propuesta.

\section{Procesamiento del lenguaje natural}

\subsection{De los sistemas de reglas a los modelos estadísticos}

% Recorrido breve: reglas y gramáticas -> métodos estadísticos -> aprendizaje
% profundo. No hace falta una historia exhaustiva, pero sí situar de dónde viene
% la tecnología que usas.

\subsection{Representación vectorial del lenguaje}

% La idea clave: representar palabras y frases como vectores para que la
% similitud de significado sea una distancia. De las bolsas de palabras a las
% incrustaciones densas \cite{reimers2019sbert}.

\subsection{La arquitectura Transformer}

% Qué aporta el mecanismo de atención y por qué es la base de todo lo que viene
% después. A nivel conceptual: no hace falta el detalle matemático.

\section{Modelos de lenguaje grandes}

\subsection{Fundamento y funcionamiento}

% Qué es un LLM y cómo se entrena, en términos comprensibles. Qué sabe hacer
% bien y qué no.

\subsection{El problema de las alucinaciones}

% La razón de ser de todo tu TFG: por qué un LLM genera información plausible
% pero falsa, y por qué eso es inaceptable en orientación académica. Ejemplos
% reales si los tienes (pregúntale a un LLM por un grado de la EPSJ y enseña lo
% que se inventa: eso vale más que tres párrafos teóricos).

\subsection{Modelos de pesos abiertos frente a modelos propietarios}

% Qué diferencia hay, qué implica cada opción (privacidad, coste, control) y por
% qué la distinción importa en un sistema que va a leer preguntas de menores.
% OJO: aquí se explica el panorama; la ELECCIÓN concreta del modelo va al Cap. 4.

\section{Recuperación de información}

\subsection{Búsqueda léxica frente a búsqueda semántica}

% Buscar por coincidencia de palabras frente a buscar por significado. Por qué
% lo segundo es lo que necesita un estudiante que no conoce la jerga.

\subsection{Incrustaciones (\textit{embeddings})}

% Qué son, qué es la ventana de entrada y por qué importa. Qué familias de
% modelos existen para el español (multilingües frente a monolingües).

\subsection{Bases de datos vectoriales}

% Qué son y en qué se diferencian de una base de datos tradicional. Panorama:
% las hay embebidas/locales y como servicio en la nube. Qué es una búsqueda de
% vecinos más cercanos aproximada y por qué hace falta.

\subsection{Métricas de evaluación de la recuperación}

% Recall@K, MRR y nDCG: qué mide cada una y qué significa un valor alto.
% IMPORTANTE: explícalas aquí para poder usarlas sin explicarlas en Resultados.

\section{Recuperación aumentada por generación (RAG)}

\subsection{Origen y motivación}

% El artículo que introduce la arquitectura \cite{lewis2020rag} y el problema
% que venía a resolver.

\subsection{Arquitectura y componentes}

% Las tres piezas: indexación, recuperación y generación. Un diagrama conceptual
% GENÉRICO de RAG encaja aquí (distinto del de tu arquitectura concreta, que va
% en el Cap. 4).

\subsection{Fragmentación de documentos}

% Por qué hay que trocear: la ventana del modelo de incrustaciones es limitada,
% un fragmento grande diluye la señal y uno pequeño pierde contexto.
% Qué estrategias existen (tamaño fijo con solape, por estructura, semántica) y
% qué asume cada una. NO digas todavía cuál elegiste: eso es el Cap. 4.

\subsection{Evaluación de sistemas RAG}

% Cómo se mide la calidad de la generación: fidelidad al contexto recuperado y
% relevancia de la respuesta. Herramientas de evaluación automática que existen.

\section{Extracción automatizada de información web}

\subsection{Fundamentos}

% Qué es el rastreo web, cómo funciona (peticiones HTTP, análisis del HTML,
% selectores) y qué problemas típicos aparecen (codificación, estructura
% inconsistente, contenido generado por JavaScript).

\subsection{Tipos de herramientas}

% Panorama: analizadores de HTML, marcos de rastreo completos y navegadores
% automatizados. Qué resuelve cada familia. La COMPARATIVA y la elección
% concreta van al Cap. 4 (ADR-0002).

\section{Trabajos relacionados}

\subsection{Asistentes conversacionales en el ámbito universitario}

% Qué han hecho otras universidades. Pruébalos de verdad si están accesibles.

\subsection{Sistemas de recomendación educativos}

% Qué enfoques existen para recomendar estudios (filtrado colaborativo, tests
% vocacionales, sistemas basados en reglas) y en qué se diferencian de un RAG.

\subsection{Posicionamiento de este trabajo}

% La sección más importante del capítulo: en qué se parece tu sistema a lo que
% existe, en qué se diferencia y qué hueco concreto llena. Aquí es donde el
% tribunal ve si el TFG aporta algo o repite lo ya hecho.

\section{Marco legal y ético}

\subsection{Protección de datos personales}

% RGPD \cite{rgpd2016} y LOPDGDD \cite{lopdgdd2018}. Qué se considera dato
% personal y qué es información anónima (Considerando 26 del RGPD). Aplicación
% al caso: la colección documental no contiene datos de personas físicas.

\subsection{Legalidad de la extracción automatizada de datos}

% Qué dice el marco jurídico sobre extraer datos publicados abiertamente por un
% organismo público, con fines académicos. Jurisprudencia de referencia
% \cite{sts572_2012}. Buenas prácticas: robots.txt, retardo entre peticiones,
% agente de usuario identificable.

\subsection{Licencias de software y de modelos}

% Licencias libres y copyleft \cite{gpl3}. Y el matiz de los modelos: pesos
% descargables no significa uso libre; hay que leer la licencia de cada modelo.

\subsection{Accesibilidad}

% Criterios WCAG \cite{wcag22} y por qué obligan a un sistema dirigido a
% estudiantes de secundaria.

\subsection{Reglamento europeo de inteligencia artificial}

% El Reglamento (UE) 2024/1689 \cite{aiact2024}: qué sistemas regula y cuál es
% la obligación de transparencia aplicable a un chatbot informativo.
